\chapter{Regularised D'Agostini iterative unfolding}

\label{ch:appendix_unfolding}

This appendix describes my understanding of the derivation and implementation of the regularised D'Agostini iterative unfolding algorithm, particularly the origin of the Equation~\ref{eq:unfolding}, based on References~\cite{DAgostini:1994fjx,DAgostini:2010hil,KuuselaThesis2016}. The reasons for the smearing of reconstructed distributions and the needs to reverse this process are discussed in Section~\ref{sec:migration_unfolding}.

The number of events of a reconstructed distribution $\boldsymbol{y}$ in a particular bin $j$ can be related to the number of events of a true distribution $\boldsymbol{\lambda}$ in the bin $\alpha$ via a migration matrix $U_{j\alpha}$ such as
\begin{equation}
    \label{eq:folding}
    y_{j} = \sum_{\alpha} U_{j\alpha}\lambda_{\alpha},
\end{equation}
where the migration matrix $U_{j\alpha}$ corresponds to
\begin{equation}
    \label{eq:migration}
    U_{j\alpha} = P(y \in \text{reco bin } j~|~\lambda \in \text{true bin}~\alpha).
\end{equation}
Here, $y$ and $\lambda$ are the events from the reconstructed distribution $\boldsymbol{y}$ and the true distribution $\boldsymbol{\lambda}$, respectively. Therefore, the Equation~\ref{eq:folding} can be also thought of as the probability that a $y$ event belongs to the bin $j$,
\begin{equation}
    \label{eq:yDistribution}
    P(y \in \text{reco bin } j) = \sum_{\alpha} P(y \in \text{reco bin } j~|~\lambda \in \text{true bin}~\alpha)  P(\lambda \in \text{true bin}~\alpha) .
\end{equation}

To reverse the smearing of the reconstructed distribution as described in Section~\ref{sec:migration_unfolding}, the unfolding matrix $\tilde{U}_{j\alpha}$ can be expressed as
\begin{equation}
    \label{eq:unfoldingMatrix}
    \tilde{U}_{j\alpha} = P(\lambda \in \text{true bin } \alpha~|~y \in \text{reco bin }j).
\end{equation}
Applying the Bayes' theorem, this becomes
\begin{equation}
    \label{eq:unfoldingMatrix2}
    \tilde{U}_{j\alpha} = \frac{P(y \in \text{reco bin } j~|~\lambda \in \text{true bin }\alpha) P(\lambda \in \text{true bin } \alpha)}{P(y \in \text{reco bin } j)}.
\end{equation}
Using the Equation~\ref{eq:yDistribution}, the denominator of the unfolding matrix can be rewritten as
\begin{equation}
    \label{eq:unfoldingMatrix3}
    \tilde{U}_{j\alpha} = \frac{P(y \in \text{reco bin } j~|~\lambda \in \text{true bin }\alpha) P(\lambda \in \text{true bin } \alpha)}{\sum_{\beta}
    P(y \in \text{reco bin } j~|~\lambda \in \text{true bin}~\beta) P(\lambda \in \text{true bin}~\beta)}.
\end{equation}

Taking into account that the conditional probability can be represented by the migration matrix from Equation~\ref{eq:migration}, the unfolding matrix is
\begin{equation}
    \label{eq:unfoldingMatrix4}
    \tilde{U}_{j\alpha} = \frac{U_{j\alpha} P(\lambda \in \text{true bin } \alpha)}{\sum_{\beta}
    U_{j\beta} P(\lambda \in \text{true bin}~\beta)}, 
\end{equation}
The zeroth order probability of the true event $\lambda$ in the bin $\alpha$ is the true event $\lambda_{\alpha}^0$ over the total of the observed events, i.e. 
\begin{equation}
    \label{eq:zerothOrder}
    P^0(\lambda \in \text{true bin}~\alpha) = \frac{\lambda_{\alpha}^0}{\sum_{j}y_{j}}, 
\end{equation}
and hence the unfolding migration matrix for the first iteration can be represented by
\begin{equation}
    \label{eq:unfoldingMatrix5}
    \tilde{U}_{j\alpha} = \frac{U_{j\alpha} \lambda_{\alpha}^0}{\sum_{\beta}
    U_{j\beta} \lambda_{\beta}^0}.
\end{equation}
As discussed in Section~\ref{sec:migration_unfolding}, the zeroth order true distribution comes from the simulation prediction of the true distribution.

The unfolded truth can be expressed as 
\begin{equation}
    \label{eq:unfolding1}
    \lambda_{\alpha} = \frac{\sum_{j} P(\lambda \in \text{true bin} ~\alpha~|~y \in \text{reco bin}~j)~y_{j}}{\sum_{i} P(y \in \text{reco bin}~i~|~\lambda \in \text{true bin}~\alpha)},
\end{equation}
where the denominator represents the overall unfolding efficiency $\epsilon_{\alpha}$, which is the sum over the reconstructed bins $i$ of the migration matrix in the Equation~\ref{eq:migration}, 
\begin{equation}
    \label{eq:efficiencyUnfolding}
    \epsilon_{\alpha} = \sum_{i} P(y \in \text{reco bin } i~|~\lambda \in \text{true bin}~\alpha) = \sum_{i} U_{i\alpha}.
\end{equation}
The estimate of the true distribution after the first iteration is then given by substituting the Equation~\ref{eq:unfoldingMatrix5} in the Equation~\ref{eq:unfolding1}, i.e. 
\begin{equation}
    \label{eq:unfolding2}
    \lambda_{\alpha} = \frac{\lambda_{\alpha}^0}{\sum_{i}U_{i\alpha}}\sum_{j} \frac{U_{j\alpha}y_j}{\sum_{\beta}U_{j\beta}\lambda_{\beta}^0}.
\end{equation}

This is valid for each iteration, as the true distribution $\boldsymbol{\lambda}$ is updated after each step with the new prediction. Finally, the $(k+1)$th iteration of the D'Agostini method can be determined by 
\begin{equation}
    \label{eq:unfolding3}
    \lambda_{\alpha}^{(k+1)} = \frac{\lambda_{\alpha}^{(k)}}{\sum_{i}U_{i\alpha}}\sum_{j} \frac{U_{j\alpha}y_j}{\sum_{\beta}U_{j\beta}\lambda_{\beta}^{(k)}}.
\end{equation}

It is also important to note that this method was shown to be equivalent to the expectation--maximisation (EM) iteration algorithm \cite{EM:477e7e2b-4ded-3369-981e-9b40850a2701} for finding the maximum likelihood estimator (MLE) of $\boldsymbol{\lambda}$ in the Poisson regression problem $\boldsymbol{y} \sim \text{Poission}(\boldsymbol{U\lambda})$, $\boldsymbol{\lambda} \geq 0$ as $k \rightarrow \infty$ \cite{McLachlan2004Algorithm, KuuselaThesis2016}.